\chapter{Related Work}
\label{ch:related}

\section{Data-Movement Characterization Methodologies}

The direct precursor to this thesis is DAMOV~\cite{oliveira2021damov}, which systematized the question: ``Which functions of a workload are data-movement-bound, and would near-data processing help?'' for CPUs. DAMOV's three-step process screen functions that matter; characterize them with hardware counters plus architecture-independent locality analysis; and classify them into bottleneck families \emph{derived from the data by clustering} is adopted here in full. Its central metric, the last-to-first miss ratio (LFMR: of the requests that miss the first cache, the fraction that also miss the last and reach DRAM), is retained in a redefined form. DAMOV's validation discipline is also followed: thresholds are frozen and stress-tested on held-out inputs, an independent clustering algorithm reproduces the classes, and cache-size sweeps confirm class behavior. Chapter~\ref{ch:method} enumerates, component by component, the GPU analogue of each DAMOV element and highlights the three places where GPUs required a redesign rather than a direct translation. Where DAMOV's third step uses a simulator-based intervention (host vs.\ host-with-prefetcher vs.\ NDP across a core-count sweep in ZSim/Ramulator), this thesis substitutes two \emph{measured} interventions on silicon independent core and memory-clock scaling, and the reuse-distance cache-capacity curve and postpones the simulated axis to generated configuration (\S\ref{sec:met-parity}).

On the GPU side, the counter-metric curation employed here follows the NCU-based roofline methodology of Cao et al.~\cite{cao2023gpudb}, who demonstrated a disciplined path from Nsight Compute counters to roofline placement for GPU database kernels. This thesis reuses their metric taxonomy (extended with stall-reason and occupancy axes) and their practice of measuring, rather than quoting, roofline ceilings. The roofline model itself is due to Williams et al.~\cite{williams2009roofline}. Reuse-distance (stack-distance) analysis originates with Mattson et al.~\cite{mattson1970stack}; its use here as a \emph{measured} hit-rate-versus-capacity curve from binary-instrumentation traces (\S\ref{sec:loc-metrics}) is the GPU translation of DAMOV's architecture-independent locality step. One novel addition, not present in the CPU literature, is filtering by memory space, without which register-spill and scratchpad traffic can corrupt the locality signal (\S\ref{sec:loc-correction}).

\section{Processing-in-Memory and Near-Data Systems}

The PiM literature spans over three decades; the modern wave is surveyed by Mutlu et al.~\cite{mutlu2019pim} and Ghose et al.~\cite{ghose2019pim}. Architectures directly relevant to this thesis include Tesseract~\cite{ahn2015tesseract} (near-memory graph processing, which serves as the lineage for scatter-capable PiM), PIM-enabled instructions~\cite{ahn2015pei} (fine-grained offload of cache-defeating operations), TOM~\cite{hsieh2016tom} (transparent offload of GPU code blocks to 3D-stacked memory, selected by compiler bandwidth-benefit analysis a methodologically close GPU-offload ancestor), and the workload-driven placement study of Boroumand et al.~\cite{boroumand2018google} for consumer workloads. Commercial near-bank silicon UPMEM DPUs~\cite{gomezluna2021upmem}, Samsung HBM-PIM~\cite{kwon2021hbmpim}, and SK~hynix GDDR6-AiM~\cite{lee2022aim} demonstrates that the substrates assigned in this thesis are realizable, not hypothetical. In-storage processing is grounded in Smart-SSD query offload~\cite{do2013smartssd} and Biscuit~\cite{gu2016biscuit}. However, none of these works provide the \emph{workload side} for Physical AI: a measured statement of which SLAM computations and data structures are best suited to each substrate. This is the gap addressed by this thesis.

\section{SLAM Systems and Their Evaluation}

The SLAM algorithmic literature is well established; Cadena et al.~\cite{cadena2016slam} provide a comprehensive survey. The systems used for accuracy validation in this thesis are cuVSLAM itself, as documented in NVIDIA's technical report~\cite{cuvslam2025}, and ORB-SLAM2/3, which serve as standard open-source references~\cite{murartal2017orbslam2,campos2021orbslam3}. Benchmarking SLAM \emph{as a computational workload} rather than solely for accuracy was pioneered by SLAMBench~\cite{nardi2015slambench}, which profiled dense SLAM across CPU and GPU platforms at whole-pipeline granularity. This thesis differs in three key ways: it measures a \emph{production} sparse V-SLAM stack rather than a research kernel; it achieves per-kernel and per-data-structure granularity; and its output is substrate placement, not a platform comparison. Trajectory-accuracy methodology follows the TUM benchmark conventions~\cite{sturm2012tum} and the evaluation tutorial by Zhang and Scaramuzza~\cite{zhang2018tutorial}, incorporating Umeyama alignment~\cite{umeyama1991}. The datasets used KITTI~\cite{geiger2013kitti}, EuRoC~\cite{burri2016euroc}, TUM RGB-D~\cite{sturm2012tum}, TUM-VI~\cite{schubert2018tumvi}, and ICL-NUIM~\cite{handa2014iclnuim} are the standard ground-truth corpora in the community.

\section{GPU Instrumentation}

Binary instrumentation of GPU machine code is provided by NVBit~\cite{villa2019nvbit}, upon which this thesis's address tracing is built (with two extensions: launch-window/kernel-name gating to limit trace volume, and an allocation-event sidecar that timestamps allocator activity using the trace's own launch-id clock). GPU simulation and power modeling Accel-Sim~\cite{khairy2020accelsim} and AccelWattch~\cite{kandiah2021accelwattch} appear in this thesis only as the targets for generated configurations in the follow-on study (\S\ref{sec:roadmap}); in accordance with project policy, no simulated numbers are reported in the results chapters.

\section{Positioning}

In summary, relative to DAMOV, this thesis is the GPU re-derivation plus the extension from bottleneck class to substrate verdict. Relative to GPU profiling practice, it adds machine-independent locality and named data-structure attribution on top of counters. Relative to the PiM systems literature, it provides the missing measured workload analysis for a deployed Physical-AI perception stack. And relative to SLAM benchmarking, it is the first study at the granularity where architectural decisions are actually made the kernel and the allocation.